\hypertarget{appendix-cbe594d1ab0a}{}\label{referencestart:appendix-cbe594d1ab0a}
\section{Integrity and authenticity}\label{appendix-cbe594d1ab0a-integrity-and-authenticity}

\hypertarget{appendix-cbe594d1ab0a-integrity-and-authenticity}{}

\textbf{Curriculum depth:} Derived with executable primitive witnesses · \textbf{First seen:} C006 · \textbf{Latest development:} Phase 28 repair

\textbf{Integrity} means detecting unauthorized or accidental modification relative to a trusted reference. \textbf{Authenticity} additionally answers who or what was authorized to create the message.

Hashing a file and storing the digest beside it protects against accidental corruption only when the digest itself is trusted. An attacker able to replace both can recompute the hash.

A message authentication code solves a different problem: parties share a secret key \(K\) and compute a tag such as

\[
\operatorname{HMAC}_K(m).
\]

Someone without the key should not be able to forge a valid tag for a modified message under the security assumptions of HMAC and the underlying hash.

A digital signature separates signing and verification keys and supports different trust models. The choice between MAC and signature is therefore architectural, not a matter of which API call is available.

\subsection{HMAC construction}\label{appendix-cbe594d1ab0a-hmac-construction}

\hypertarget{appendix-cbe594d1ab0a-hmac-construction}{}

For hash block size \(B\), normalize key \(K\) to one block, then compute

\[
\operatorname{HMAC}_K(m)
=
H((K\oplus opad)\,\|\,
  H((K\oplus ipad)\,\|\,m)).
\]

The inner and outer domains prevent the construction from being equivalent to a naive hash of key concatenated with message.

tools/labs/crypto\_primitives\_lab.cpp checks HMAC-SHA256 against an RFC 4231 known-answer vector.

\subsection{Key management is part of authenticity}\label{appendix-cbe594d1ab0a-key-management-is-part-of-authenticity}

\hypertarget{appendix-cbe594d1ab0a-key-management-is-part-of-authenticity}{}

A mathematically correct MAC does not help if:

\begin{itemize}
\tightlist
\item
  the key is embedded publicly in the executable;
\item
  keys are reused across incompatible protocols without domain separation;
\item
  a compromised key is never rotated;
\item
  comparison leaks useful timing information in an adversarial environment;
\item
  replayed old authenticated messages are accepted without nonce/counter/version policy.
\end{itemize}

Cryptographic primitives are building blocks, not complete protocols.

\subsection{Production boundary}\label{appendix-cbe594d1ab0a-production-boundary}

\hypertarget{appendix-cbe594d1ab0a-production-boundary}{}

The educational implementations in this repository exist to make the bytes/rounds/contracts reconstructable. Production systems should use vetted cryptographic libraries and platform key-storage facilities rather than treating a teaching SHA/HMAC implementation as hardened security code.

Related: \hyperref[appendix-ed338677de0a]{Checksums CRCs and hashes}, \hyperref[appendix-50af792c9da7]{HMAC and HKDF}, \hyperref[appendix-098192f8d46d]{Merkle trees and authenticated structure}, \hyperref[appendix-5888fd6102dc]{Shamir secret sharing}.
